\section{未来系统方向：hybrid serving、test-time infrastructure 与 continual learning}

\videofigure{figures/research-directions.jpg}{未来研究覆盖 synthetic data 的学习原则、continual learning，以及高吞吐低延迟的 test-time scaling 基础设施。}{00:52:18--00:57:06}

Test-time scaling 不再是一次 forward pass：它包含重复采样、多轮 agent 更新、tool call、KV cache 重用、并行分支和 verifier。传统 chatbot serving 针对单轮解码优化，无法自动满足这种不规则 workload。

\begin{importantbox}{推理基础设施必须感知 agent graph}
Scheduler 需要看到分支依赖、可共享 prefix、工具等待和候选优先级，才能在延迟、吞吐与能耗之间做全局优化。
\end{importantbox}

\videofigure{figures/intelligence-efficiency-future.jpg}{下一阶段需要 local/cloud 混合 serving、新模型与 kernel，以及更细粒度的 intelligence-per-watt 测量。}{00:57:06--00:59:02}

可把动态路由写成成本敏感决策：

$$
a^*=\arg\min_{a\in\{\text{local},\text{cloud}\}}
\left(C_a+\lambda_L L_a+\lambda_E E_a+\lambda_R R_a\right),
$$

\begin{itemize}
    \item $C_a$：货币成本；
    \item $L_a$：延迟；
    \item $E_a$：能耗；
    \item $R_a$：失败、隐私或安全风险；
    \item $\lambda$：应用对各项代价的权重。
\end{itemize}

\subsection{Memory 与 weight update 的分工}

课程问答讨论 continual learning：更新外部 memory 更容易、可追踪，也不易破坏已有能力；但若目标是获得新的 reasoning skill、跨领域抽象或机器人跨 embodiment 泛化，仅增加 memory 未必足够，仍可能需要参数更新。

\begin{knowledgebox}{Knowledge acquisition 与 skill acquisition 不同}
数据库适合存事实、历史和用户偏好；weight update 更适合改变模型如何组合信息、如何行动和如何迁移技能。实际系统需要决定哪些经验写 memory，哪些经验进入训练。
\end{knowledgebox}

\begin{warningbox}{无限 context 不是免费 continual learning}
即使能存下所有经历，模型也可能无法在超长 context 中准确检索和组合。有效 memory 仍需压缩、索引、遗忘策略与冲突处理。
\end{warningbox}

\subsection{本章小结}

未来 agent 系统将同时优化模型、serving、硬件和学习机制：本地与云端动态路由，基础设施支持不规则 test-time search，memory 保存可检索经验，weight update 学习可迁移技能。

